%%%%%%%%%%%%%%%%%%%%%%%%%%%%%%%%%%%%%%%%%%%%%%%%%%%%%%%%%%%
%%%
%%% neue, eigene Kommandos
%%%
%%%%%%%%%%%%%%%%%%%%%%%%%%%%%%%%%%%%%%%%%%%%%%%%%%%%%%%%%%%

% normaler Text
% ==============
%unterstreichen
\newcommand{\ul}[1]{\underline{#1}}

%für dicke Symbole
\newcommand{\bs}{\boldsymbol}

%Gänsefüsschen unten/oben
\newcommand{\gfu}{"`}
\newcommand{\gfo}{"' }
\newcommand{\gf}{}


% Bilder -- Größe des Bildes und der
% Achsbeschriftungen wählen.
% ========================================
\newcommand{\BildSubfigureBreite}{0.45\textwidth}
\newcommand{\BildSubfigureBreiteDrittel}{0.3\textwidth}

\newcommand{\BildBreiteGross}{1\textwidth}
\newcommand{\BildBreiteMittel}{0.8\textwidth}
\newcommand{\BildBreiteKlein}{0.65\textwidth}

% caption in subfigures
\newcommand{\captionsf}[1]{\emph{#1}}



% Einheiten
% alle Einheiten müssen mit der Umgebung $...$ versehen
% werden.
% =======================================================
% Logik:
% \einh         : einfache Einheit, z.~B. m
% \einhbruch    : Einheit mit einem Bruch, z.~B. m/s
% \einhmit      : Einheit in Klammern, z.~B.  t \einhmit{s}
% \einhmitbruch : Einheit in Klammern mit Bruch, z.~B. [m/s]
%----------------------------------------------------------
%Schriftgröße der Einheiten hier vorgeben:
% \es = Kommando für die Einheiten-Schriftgröße
% mögl.: \tiny \scriptsize \footnotesize \small \normalsize
\newcommand{\es}{\normalsize}

% \einh         : einfache Einheit, z.~B. m
\newcommand{\einh}[1]{\mbox{\es{\,$\unit{#1}$}}}

% \einhbruch    : Einheit mit einem Bruch, z.~B. m/s
\newcommand{\einhbruch}[2]{\mbox{\es{\,$\unitfrac{#1}{#2}$}}}

% \einhmit      : Einheit in Klammern, z.~B.  t \einhmit{s}
\newcommand{\einhmit}[1]{\mbox{\es{$\  [\unit{#1}$]}}}

% \einhmitbruch : Einheit in Klammern mit Bruch, z.~B. [m/s]
\newcommand{\einhmitbruch}[2]{\mbox{\es{$\  [\unitfrac{#1}{#2}]$}}}


% Tabellen
% =========
% Zeilenabstände in Tabellen
\newcommand{\abstand}{\rule{0in}{.6ex} \\[.2ex]}


% Aufzählungen
% =============
%Aufzählung in einer Tabelle (kein Abstand zum
%vorangehenden Text)
\newenvironment{AufzTabelle}{
    \begin{list}{$\bullet$}{
    \setlength{\leftmargin}{0,5cm}
    \setlength{\topsep}{0cm}
    \setlength{\parskip}{0cm}
    \setlength{\partopsep}{0cm}
    \setlength{\parsep}{0 cm}
    \setlength{\itemsep}{0cm}
    }}
    {
    \end{list}
    }

% Aufzählung benutzerdefiniert
% Aufruf: \begin{Aufzaehlung}{$\bullet$}
\newenvironment{Aufzaehlung}[1]{
    \begin{list}{#1}
    {
    \setlength{\leftmargin}{2.5cm}
    \setlength{\rightmargin}{0,5cm}
    \setlength{\labelsep}{0,5cm}
    \setlength{\labelwidth}{0,5cm}
    \setlength{\topsep}{0,5cm}
    \setlength{\parskip}{0cm}
    \setlength{\partopsep}{0cm}
    \setlength{\parsep}{0cm}
    \setlength{\itemsep}{0,2cm}
    \setlength{\listparindent}{0cm}
    }
    }
    {\end{list}}

% Aufzählung benutzerdefiniert - mit Zähler
% Aufruf: \begin{AufzaehlungZaehler}{Schritt}{:}
%\newenvironment{AufzaehlungZaehler}[2]{
\newcounter{zaehler}
\newenvironment{AufzaehlungZaehler}[2]{
    \setcounter{zaehler}{0}
    \begin{list}{#1 \thezaehler #2}
    {
    \usecounter{zaehler}
    \setlength{\leftmargin}{3cm}
    \setlength{\rightmargin}{0,5cm}
    \setlength{\labelsep}{0,5cm}
    \setlength{\labelwidth}{2.5cm}
    \setlength{\topsep}{0,5cm}
    \setlength{\parskip}{0cm}
    \setlength{\partopsep}{0cm}
    \setlength{\parsep}{0cm}
    \setlength{\itemsep}{0,2cm}
    \setlength{\listparindent}{0cm}
    }
    }
    {\end{list}}

% Formeln > Abkürzungen
% ======================
\newcommand{\psip}{\dot{\psi}}
\newcommand{\w}{\omega}
\newcommand{\pfeil}{\longrightarrow}


% Formeln
% ========
%für Matrizen
\newcommand{\M}[1]{{\textbf{#1}}}

%für Vektoren in Koordinatenschreibweise (=Spaltenmatrizen)
\newcommand{\V}[1]{{\underline{#1}}}

% Vektoren mit einzelnen Elementen mit eckigen Klammern
\newcommand{\Vektor}[1]{\begin{bmatrix} #1 \end{bmatrix}}

%für komplexe Zahlen
\newcommand{\K}[1]{{\overline{#1}}}

%komplexe Vektoren
\newcommand{\VK}[1]{\K{\V{#1}}}

%komplexe Matrizen
\newcommand{\KM}[1]{\overline{\bf#1}}

%kleine Boxen in Formeln, mehrzeilig
\newcommand{\txt}[1]{
 \begin{minipage}{1.5 cm}
 \vspace{1mm}
 \begin{center} #1  \end{center}
 \end{minipage}}

%kleine Boxen in Formeln, gedrehter Text!
\newcommand{\drehtxt}[1]{
 \begin{minipage}{0.52 cm}
 \vspace{10mm}
 \begin{center}
 \begin{rotate}{90}
 \tiny
 #1
 \end{rotate}
 \end{center}
 \end{minipage}}

\newcommand{\mf}[1]{{\mathbf{#1}}}

% Datum Festlegen
\newcommand{\OrtDatum}{Lippstadt, 25. Mai 2007}

% Umlaute in Listings
\lstset{literate=%
  {Ö}{{\"O}}1
  {Ä}{{\"A}}1
  {Ü}{{\"U}}1
  {ß}{{\ss}}1
  {ü}{{\"u}}1
  {ä}{{\"a}}1
  {ö}{{\"o}}1
}

% Abbildung heißt kurz Bild
\renewcaptionname{ngerman}{\figurename}{Bild}